\section{Introduction}

Surface complexation is a key control on reactive transport in porous media,
but most pore-scale implementations assume either explicitly resolved grain
surfaces or simplified source terms that are not tied cleanly to a
conservative transport balance. The objective of this work is to develop a DBS
formulation that couples PHREEQC chemistry to unresolved internal surface area
while preserving mathematically correct inventory balances.

The guiding constraints of this implementation are:
\begin{itemize}
\item the transport equation must conserve mobile aqueous inventory in the
physical mobile-water fraction,
\item local chemistry must conserve total inventory while transferring mass
between aqueous and adsorbed states,
\item the baseline method must not rely on concentration clipping, damping, or
heuristic face-based reaction reconstruction.
\end{itemize}

The solver developed in this work is \texttt{reactiveScalarTransportDBSFoam}.
It is built on a staged validation program so every major reduction path
remains testable. This manuscript follows the same overall structure as the
earlier GeoChemFoam reactive transport paper: introduction, model description,
implementation, validation, and then progressively more complex test cases.
